% Generated draft: requires mathematical and source review.
\documentclass[11pt]{article}
\usepackage[margin=25mm]{geometry}
\usepackage{fontspec}
\setmainfont{Latin Modern Roman}
\usepackage{amsmath,amssymb,mathtools,unicode-math}
\setmathfont{Latin Modern Math}
\usepackage{enumitem,needspace,xurl,hyperref}
\setlength{\parindent}{0pt}
\setlength{\parskip}{5pt}
\setlength{\emergencystretch}{3em}
\newcommand{\N}{\mathbb{N}}
\newcommand{\Z}{\mathbb{Z}}
\newcommand{\Q}{\mathbb{Q}}
\newcommand{\R}{\mathbb{R}}
\newcommand{\C}{\mathbb{C}}
\newcommand{\F}{\mathbb{F}}
\newcommand{\A}{\mathbb{A}}
\newcommand{\PP}{\mathbb{P}}
\newcommand{\E}{\mathbb{E}}
\newcommand{\eps}{\varepsilon}
\newcommand{\dd}{\,\mathrm{d}}
\begin{document}
\section*{UnsolvedMath ID: 11000061}
\textbf{Source problem code:} catalogue entry
\subsection{Definitions and mathematical statement}
Let \(\mathcal I_g\) be the Torelli group and let \(K\) be either \(\Q\) or \(\F_2\). Define the graded subalgebra generated by degree one as
\[A_K^*=\langle H^1(\mathcal I_g;K)\rangle\subseteq H^*(\mathcal I_g;K),\qquad
A_K^n=\operatorname{span}_K\{x_1\smile\cdots\smile x_n:x_i\in H^1(\mathcal I_g;K)\},\]
with the unit in degree zero. The problem is to determine these subalgebras, in all degrees, for both fields, including the image and relations of cup products. For \(K=\Q\), the cup map factors through the exterior algebra on \(H^1\). For \(K=\F_2\), degree-one squares need not vanish, and the corresponding free commutative source is a symmetric algebra rather than an exterior algebra unless extra square relations are proved. The problem does not assert that \(A_K^*=H^*(\mathcal I_g;K)\).

The source states that \(H^*(\mathcal I_g;K)\) is a module over \(\operatorname{Sp}(2g,K)\). Conjugation by mapping classes gives an action through the integral symplectic quotient; extension to the rational algebraic group or descent to \(\operatorname{Sp}(2g,\F_2)\) requires the precise representation conventions, not just the notation. This supplied structural assertion is preserved subject to that clarification.

For rational coefficients the source reports a degree-two solution by Hain [Ha3], and a degree-three solution by Sakasai [Sa] up to one unknown piece, using symplectic representation theory over \(\R\) to simplify computations. The unresolved piece and genus ranges are not specified. For \(\F_2\), it emphasizes that needed representation facts may fail or lie beyond available modular representation methods, and reports geometric abelian-cycle calculations in [BF a2]. Specifically it gives the pullbacks
\[\sigma^*:H^2(B_3;\F_2)\to H^2(\mathcal I_g;\F_2),\]
\[(\sigma|_{\mathcal K_g})^*:H^2(B_2;\F_2)\to H^2(\mathcal K_g;\F_2),\]
and says each image has dimension ``at least \(O(g^4)\).'' The groups \(B_2,B_3\), homomorphism \(\sigma\) and precise Johnson-kernel \(\mathcal K_g\) conventions are not defined here. Big-\(O\) denotes an upper-bound class, so that phrase does not literally specify a lower bound; an intended \(\Omega(g^4)\) statement would require constants, range and a verified reading. No corrected quantitative theorem is silently inserted.
\subsection{Short English statement}
Determine all cup products generated by first cohomology of the Torelli group over \(\Q\) and \(\F_2\). The source reports low-degree progress, including an unresolved rational degree-three piece and mod-two lower bounds whose notation needs correction.
\subsection{Sources}
\begin{itemize}
\item Local UnsolvedMath registry, permanent ID \(11000061\), exact source code \texttt{catalogue entry}.
\item Farb, \emph{Problems on Mapping Class Groups and Related Topics}, Problem 5.16, PDF page 49: \url{https://www.math.uchicago.edu/~farb/papers/mcgbook.pdf}. Registry access 2026-07-29; pdf-font-structure extraction.
\item Hain [Ha3], Sakasai [Sa], and [BF a2] are cited for the respective low-degree and abelian-cycle results; full bibliographic entries are absent. The summary also mentions a Malcev approach without an exact reference.
\end{itemize}
\end{document}
